\documentclass[11pt]{article}

\usepackage[margin=1in]{geometry}
\usepackage{amsmath,amssymb}
\usepackage{booktabs}
\usepackage{graphicx}
\usepackage{float}
\usepackage{hyperref}
\usepackage[numbers,sort&compress]{natbib}

\floatplacement{figure}{H}

\newcommand{\resultsdir}{figures}
\newcommand{\vect}[1]{\boldsymbol{#1}}

\title{Physics-Informed Neural Networks for the One-Dimensional\\
Convection--Diffusion Equation}
\author{Elijah Kin}
\date{\today}

\begin{document}
\maketitle

\begin{abstract}
  We compare physics-informed neural networks (PINNs) with a limited
  Lax--Wendroff finite-volume scheme for the one-dimensional linear and
  nonlinear convection--diffusion equation with a discontinuous initial
  condition, an oscillatory Dirichlet inflow, and a zero-gradient outflow. A
  three-grid Richardson extrapolation of the limited Lax--Wendroff scheme
  provides the reference. Its $L^1$ self-convergence estimate is
  $1.3$--$5.6\times10^{-5}$ for $\nu\ge0.002$ and
  $1.4$--$3.9\times10^{-4}$ at the stress-test value $\nu=10^{-4}$; the
  corresponding $L^\infty$ estimate on the common evaluation grid is
  $5.7\times10^{-4}$--$5.4\times10^{-3}$ and $4\times10^{-2}$--$10^{-1}$,
  respectively. Across six cases, the limited scheme has four to nine times
  lower $L^1$ error than a baseline $\tanh$ PINN at $\Delta x=\pi/200$. The
  PINN's $L^\infty$ error is dominated by the initial datum: a hard branch
  that imposes the datum exactly removes this contribution without adding
  parameters. Across the exact-IC sweep of eight plain architectures (five
  seeds each), the mean $L^\infty$ is $0.25$--$0.43$ on the linear problem and
  $0.18$--$0.35$ on the nonlinear problem, compared with $0.70$--$0.80$ for the
  soft-IC $4\times64$ baseline. The lowest mean PINN values ($0.25$--$0.27$
  linear, $0.18$--$0.19$ nonlinear) remain about $2.5$--$3$ times the
  finite-volume values on the same grid. Exact initial conditions therefore
  close neither the $L^1$ nor the $L^\infty$ gap. The lowest observed mean
  linear $L^1$ is the $4\times64$ sine network; larger or deeper sine networks
  do not improve the tradeoff. The lowest observed mean nonlinear $L^\infty$
  is the $4\times64$ or $6\times64$ $\tanh$ network; differences among the best
  architectures are within one standard deviation in both cases, so these
  rankings are indicative rather than statistically resolved. All PINN runs
  use 3000 epochs; a baseline run takes about $2.4$ min on a single NVIDIA
  GeForce RTX 3060, and the hard-IC sweep was run with multiple jobs
  concurrently on the same GPU.
\end{abstract}

\section{Background: the convection--diffusion equation}
\label{sec:background}
The convection--diffusion equation couples advection, which transports $u$ with
speed $f'(u)$, with diffusion, which smooths gradients at rate $\nu$. For
$f=au$ a step becomes a heat front of thickness $O(\sqrt{\nu t})$; for
$f=u^2/2$ (Burgers' equation) a compressive step steepens into a shock of width
$O(\nu/\Delta u)$, whereas the expansive step used here opens into a
rarefaction, and the oscillatory inflow generates shocks downstream. Thin
fronts are hard for low-order finite-volume
schemes unless limited \citep{leveque2002}, and for smooth $\tanh$ networks,
which are spectrally biased toward low frequencies \citep{rahaman2019} and need
explicit high-frequency features to fit a moving discontinuity
\citep{tancik2020}. The convection-dominated limit is difficult for PINNs. We
ask how a baseline PINN compares with a limited finite-volume scheme on this
problem, whether exact initial conditions close the $L^\infty$ gap, and which
architecture is best once the initial layer is removed.

\section{Problem formulation}
\label{sec:problem}
We solve
\begin{equation}
  u_t + f_x = \nu u_{xx}, \qquad 0 \le x \le \pi,\quad t \ge 0,
  \qquad
  f =
  \begin{cases} u & \text{linear},\\ u^2/2 & \text{nonlinear},
  \end{cases}
\end{equation}
with the discontinuous initial condition
\begin{equation}
  u(x,0) =
  \begin{cases} 1/2 & x < \pi/2,\\ 3/2 & x \ge \pi/2,
  \end{cases}
  \label{eq:ic}
\end{equation}
the oscillatory Dirichlet inflow condition
\begin{equation}
  u(0,t) = 1 - \tfrac14\cos(4t) - \tfrac14\cos(8t),
\end{equation}
and the zero-gradient outflow condition $u_x(\pi,t) = 0$. We consider
$\nu \in \{0.012, 0.004, 0.002\}$ and, as a stress test, $\nu = 10^{-4}$.
For $\nu = 0$ the linear problem has the closed-form traveling solution used
to validate the classical solvers.

With unit advection speed the final time $t=\pi$ is one flow-through. The
initial jump leaves the domain before $t=\pi$, but its diffusive imprint
persists and remains difficult for the network to represent. The inflow
supplies a continuously varying signal at frequencies $\omega\in\{4,8\}$, and
the solution contains a smooth traveling pulse together with a sharp, decaying
front.

\section{Methods}
\label{sec:methods}

\subsection{Classical finite-volume scheme and reference}
The comparison solver combines a limited Lax--Wendroff convective flux with a
trapezoidal (Crank--Nicolson) diffusion term and solves a tridiagonal system
for the interior unknowns. It uses the van Leer monotonized-central (MC)
limiter. The left boundary is Dirichlet; the right boundary uses the
zero-gradient condition $u_{j_{\max}} = u_{j_{\max}-1}$. All finite-volume and
reference solves use the CFL number $\Delta t/\Delta x=0.5$.

The limiter is applied to the convective interface flux in the standard
limited Lax--Wendroff form \citep{leveque2002}; the diffusion terms are
centered in time and space.

The reference is the three-grid Richardson extrapolation of the limited
Lax--Wendroff (van Leer MC) scheme. Because the limited scheme is only
first order across the moving discontinuity, a two-grid combination with $p=2$
leaves the leading $O(\Delta x)$ front term. We combine
$\Delta x=\pi/400$, $\pi/800$ and $\pi/1600$,
\begin{equation}
  u_{\mathrm{ref}} = u_1 + 2\,(u_1-u_2) - \tfrac13\,(u_1-u_4),
  \label{eq:rich3}
\end{equation}
with $u_1$ the $\pi/1600$ solution and $u_2$, $u_4$ the $\pi/800$ and $\pi/400$
fields interpolated onto the fine grid. This cancels both the $O(\Delta x)$ and
$O(\Delta x^2)$ terms in an averaged sense; the cancellation is not pointwise
at the front. The combination is applied only for $t>0$; at $t=0$ the reference
is set to the exact initial datum, because extrapolating across the
discontinuity produces nonphysical pointwise values.

Repeating the combination one grid finer changes the field by an $L^1$ of
$1.3$--$1.6\times10^{-5}$ at $\nu=0.012$, $1.4$--$3.0\times10^{-5}$ at
$\nu=0.004$, and up to $5.6\times10^{-5}$ at $\nu=0.002$; at $\nu=10^{-4}$ it
is $1.4$--$3.9\times10^{-4}$. On the common $201\times101$ evaluation grid the
corresponding $L^\infty$ changes are $5.7\times10^{-4}$--$1.5\times10^{-3}$ for
$\nu\ge0.004$, up to $5.4\times10^{-3}$ at $\nu=0.002$, and
$3.9\times10^{-2}$--$1.0\times10^{-1}$ at $\nu=10^{-4}$. The $L^\infty$ error is
localized at the front and is much larger than the $L^1$ error. We take the
repeated-combination differences as the self-convergence estimate.

\subsection{Physics-informed neural network}
\label{sec:pinn}
The surrogate is a physics-informed neural network \citep{raissi2019}: a
fully connected feed-forward network $u_\theta:\mathbb{R}^2\to\mathbb{R}$ that
takes the raw space-time coordinates $(t,x)$ as a two-dimensional input and
returns the solution at that point. The baseline
architecture is an input projection to width $H=64$, four hidden-to-hidden
layers of width $64$, and a linear output, with $\tanh$ activations and
$16{,}897$ trainable parameters; we label this network $4\times64$. The input
is not rescaled; the network
sees the physical coordinates in $[0,\pi]$, so no coordinate features are
supplied.

The parameters are found by minimizing the composite loss
\begin{equation}
  \mathcal{L} = \lambda_{\text{pde}}\mathcal{L}_{\text{pde}}
  + \lambda_{\text{ic}}\mathcal{L}_{\text{ic}}
  + \lambda_{\text{bc}}\mathcal{L}_{\text{bc}},
  \label{eq:loss}
\end{equation}
where $\mathcal{L}_{\text{pde}}$ is the mean-squared PDE residual
$u_t + f_x - \nu u_{xx}$ evaluated by automatic differentiation at interior
collocation points, and $\mathcal{L}_{\text{ic}}$, $\mathcal{L}_{\text{bc}}$
are mean-squared penalties enforcing the initial and boundary data
(including the zero-gradient outflow condition). The residual is evaluated on
$N_{\text{pde}}=3000$ uniformly sampled interior points, with
$N_{\text{ic}}=128$ initial samples and $N_{\text{bc}}=128$ samples per
boundary.
Each epoch shuffles the collocation points and takes mini-batches of size
$256$. We use Adam with learning rate $10^{-3}$, decayed by $0.5$ at epochs
$1200$ and $2200$, and fixed weights
$(\lambda_{\text{pde}},\lambda_{\text{ic}},\lambda_{\text{bc}}) = (1,10,1)$.
All PINN results below are trained for 3000 epochs on a single NVIDIA
GeForce RTX 3060 (12\,GB). The hard-IC sweep was executed with multiple jobs
concurrently on the same GPU, so its per-run wall-clock times reflect GPU
contention rather than the cost of the exact-IC branch, which adds negligible
work.

\subsection{Architecture variants}
\label{sec:variants}
We consider two activations, $\tanh$ and sine, at four sizes:
$2\times32$, $4\times64$ (baseline), $6\times64$, and $4\times128$, spanning
$2{,}241$ to $66{,}561$ parameters. The sine networks use $\sin(3z)$ and the
SIREN initialization of Sitzmann et al. \citep{sitzmann2020}. This gives eight
architectures; each is trained from five seeds on the linear and nonlinear
problems at $\nu=0.012$ under the common 3000-epoch budget.

\subsection{Exact initial condition}
\label{sec:hardic}
All architectures in Section~\ref{sec:variants} use the same treatment of the
discontinuous datum. Let $u_\theta$ denote the raw network output. The model
used for training and evaluation is
\begin{equation}
  \tilde u_\theta(t,x) =
  \begin{cases}
    u_0(x), & t=0,\\
    u_\theta(t,x), & t>0,
  \end{cases}
  \label{eq:hardic}
\end{equation}
where $u_0$ is the step of Eq.~\eqref{eq:ic}. At the $t=0$ collocation and
evaluation points the output is the exact datum, independent of the network
parameters; the branch adds no parameters. The PDE and boundary terms of
Eq.~\eqref{eq:loss} are evaluated on $\tilde u_\theta$; the IC term is
evaluated on the raw field $u_\theta$, because $\tilde u_\theta$ is identically
$u_0$ at $t=0$ and would give no gradient. The penalty keeps the $t>0$ branch
close to the datum as $t\downarrow0$, so the branch does not introduce a large
inconsistency.

This is a hard-constraint construction in the spirit of
\citep{lagaris1998,lu2021hpinn,hao2024hardic}, applied here to a discontinuous
datum. Rather than the usual $u_0(x)+tN_\theta(x,t)$ ansatz, which suppresses
the network correction near $t=0$, we switch branches at $t=0$ and train the
$t>0$ branch with the IC penalty. The soft-constraint baseline uses the same
architectures without the branch and with the IC penalty applied directly to
the network output.

\section{Results}
\label{sec:results}

Errors are reported as $L^1$ (mean absolute) and $L^\infty$ (maximum) on a
common $201\times101$ space-time grid against this reference; the
finite-volume fields are sampled at the same time levels as the PINN. Because
the maximum is sensitive to the time sampling, all methods are scored on this
common grid. The finite-volume scheme starts from the exact discontinuous
datum, so its error vanishes at $t=0$, whereas the soft-IC PINN must learn the
initial condition and its largest error is there. These norms are used only
for evaluation; training minimizes the composite loss of Eq.~\eqref{eq:loss}.

\subsection{Accuracy over the six cases}
\label{sec:accuracy}
The finite-volume scheme is consistently more accurate than the baseline
PINN. At $\Delta x=\pi/200$ its $L^1$ error on the common grid is
$(1.6$--$2.1)\times10^{-3}$, compared with $(0.8$--$1.6)\times10^{-2}$ for the
PINN, so the scheme has four to nine times lower error; the largest ratio is
$9.4$ (linear, $\nu=0.004$) and the smallest $4.2$ (nonlinear, $\nu=0.012$).
The gap is similar for both problems. PINN values are five-seed means at the
3000-epoch budget, scored against the same reference as the finite-volume
solutions.

The soft-IC baseline's $L^\infty$ error is $0.70$--$0.80$. Most of it is at
$t=0$: removing that frame lowers the seed-averaged error to $0.22$--$0.47$,
still above the finite-volume $0.07$--$0.20$ on the same grid. Some individual
runs remain as high as $0.54$ even after removing that frame, so the initial
layer explains the average but not every run. The $L^1$ error behaves
differently: removing $t=0$ leaves it nearly unchanged, consistent with the
initial layer being a small fraction of the space-time average; the exact-IC
results below confirm that exact initial conditions do not close the $L^1$
gap. With soft constraints, architecture changes help only unevenly; a
sine activation brings the linear $\nu=0.012$ error to $2.3$ times the
finite-volume value but does not improve the nonlinear problems.
Section~\ref{sec:arch} shows that exact initial conditions reduce the
$L^\infty$ gap, but the finite-volume scheme remains more accurate.

\subsection{Linear convection--diffusion}
Figure~\ref{fig:linear} shows the final-time profile of the exact-IC sine
network; at this time it tracks the reference closely, with the visible
deviation near the outflow. The limited scheme is second order in smooth
regions but only first order across the moving front, so whole-grid $L^1$
convergence orders are about $1.0$--$1.4$. At $t=0$ the soft-IC PINN's
$L^\infty$ error is at $x=\pi/2$. An optimally smeared step would give a
half-jump error of $0.5$; here the error is larger because the learned
transition is displaced as well as smoothed: at $x=\pi/2$ the network returns
$0.69$ where the Richardson-extrapolated solution is $1.5$, an error of
$0.81$. As the solution evolves, the maximum error over the domain decays
(from $0.81$ at $t=0$ to $0.27$ at $t=0.5$ and $0.08$ at $t=2$), consistent
with an error that originates at the initial datum. The exact initial
condition of Section~\ref{sec:hardic} removes this contribution; the residual
whole-grid $L^\infty$ occurs shortly after $t=0$ at the front, while at the
final time the profile is nearly indistinguishable from the reference except
near the outflow.

\begin{figure}
  \centering
  \includegraphics[width=0.62\textwidth]{\resultsdir/snapshot_linear_nu0.012_hard_ic.png}
  \caption{Linear $\nu=0.012$: final-time profiles for the exact-IC
    $4\times64$ sine network (seed 2). The solid black curve is the Richardson-extrapolated
    Lax--Wendroff reference, the red dashed curve is the PINN, and the
    dotted blue curve is the limited Lax--Wendroff solution at
  $\Delta x=\pi/200$.}
  \label{fig:linear}
\end{figure}

\subsection{Low-diffusivity stress test}
As $\nu$ decreases the front sharpens. At $\nu=10^{-4}$ the exact-IC
$4\times64$ sine model (linear) and the $4\times64$ $\tanh$ model (nonlinear),
each with three seeds, have $L^\infty=0.68\pm0.19$ and $0.63\pm0.05$, compared
with $0.34$--$0.41$ for the limited scheme on the same grid; their $L^1$ errors
are $8.3\times10^{-3}$ and $4.1\times10^{-2}$, versus $2.4\times10^{-3}$ and
$3.9\times10^{-3}$ (Figure~\ref{fig:stress}). The exact IC removes the
initial-layer contribution (the soft-IC baseline had
$L^\infty\approx0.87$--$0.90$), but the moving front and, in the nonlinear
case, the downstream shocks now dominate; in the linear case the largest
errors for two of the three seeds are at the outflow.

The reference $L^\infty$ self-convergence on the common grid at this value is
$3.9\times10^{-2}$ (linear) and $1.0\times10^{-1}$ (nonlinear), so the
stress-test $L^\infty$ values carry a reference uncertainty of roughly
$10$--$25\%$ for the finite-volume scheme and $6$--$16\%$ for the PINN. The
ordering is unchanged, but the absolute values are less precise than at
$\nu=0.012$. The classical $L^1$ continues to decrease under refinement (to
$9.0\times10^{-4}$ and $1.9\times10^{-3}$ at $\pi/400$). The nonlinear
$L^\infty$ at the unresolved viscous shock stays near $0.35$--$0.41$, as
expected while the shock remains under-resolved; the linear $L^\infty$ falls
to $0.19$ at $\pi/400$ as the diffusive front is resolved.

\begin{figure}
  \centering
  \includegraphics[width=0.48\textwidth]{\resultsdir/analysis_stress_linear_hard_ic.png}
  \includegraphics[width=0.48\textwidth]{\resultsdir/analysis_stress_nonlinear_hard_ic.png}
  \caption{Stress test at $\nu=10^{-4}$: linear (left) and nonlinear (right).
    The solid black curve is the Richardson-extrapolated Lax--Wendroff
    reference, the dotted blue curve is the limited Lax--Wendroff solution at
    $\Delta x=\pi/200$, and the red dashed curve is the exact-IC PINN
  ($4\times64$ sine for linear, $4\times64$ $\tanh$ for nonlinear).}
  \label{fig:stress}
\end{figure}

\subsection{Exact initial condition and architecture}
\label{sec:arch}
Table~\ref{tab:arch} compares the eight plain architectures of
Section~\ref{sec:variants} with exact initial conditions, each trained from
five seeds at $\nu=0.012$. The branch adds no parameters, so the exact-IC rows
are architecture-for-architecture comparable. The table includes soft-IC
controls only for the two $4\times64$ networks; for those, exact IC removes the
initial-layer error, and across the full exact-IC sweep the mean $L^\infty$ is
$0.25$--$0.43$ on the linear problem and $0.18$--$0.35$ on the nonlinear
problem, compared with $0.70$--$0.80$ for the soft-IC $4\times64$ baseline. The
lowest means remain above the finite-volume values ($0.09$
linear, $0.07$ nonlinear): on the nonlinear problem the $4\times64$ and
$6\times64$ $\tanh$ networks give the lowest $L^\infty$, while the $4\times64$
sine network gives the lowest $L^1$. On the linear problem the $4\times128$
sine network gives the lowest $L^\infty$ and the $4\times64$ sine network the
lowest $L^1$. Differences among the best few architectures are within one
standard deviation, so these rankings are indicative rather than
statistically resolved.

Larger or deeper sine networks do not improve the linear tradeoff. The
$6\times64$ sine network becomes unstable on the nonlinear problem, where one
seed reaches $L^1=0.207$, and the $4\times128$ sine network gives a higher
$L^1$ than the $4\times64$ sine network on both problems. The $4\times64$ sine
network is therefore the better sine choice for $L^1$, while the $4\times128$
sine network buys only a small linear $L^\infty$ gain.

Exact IC is primarily an $L^\infty$ correction and closes neither gap: the
best $L^1$ and $L^\infty$ values remain about $2.5$--$3$ times those of the
finite-volume scheme. On the linear problem it slightly worsens $L^1$ relative
to the soft-IC sine baseline, while on the nonlinear problem it improves it;
both differences are within one standard deviation. Once the initial layer is
removed, architecture choice sets the remaining
accuracy, but the finite-volume scheme remains the more accurate method.

\begin{table}[t]
  \centering
  \small
  \setlength{\tabcolsep}{4pt}
  \begin{tabular}{llcccc}
    \toprule
    Variant & Setting & \multicolumn{2}{c}{Linear} & \multicolumn{2}{c}{Nonlinear}\\
    \cmidrule(lr){3-4}\cmidrule(lr){5-6}
    & & $L^1\,[10^{-3}]$ & $L^\infty$ & $L^1\,[10^{-3}]$ & $L^\infty$\\
    \midrule
    \multicolumn{6}{l}{\emph{Soft IC (5 seeds, 3000 epochs)}}\\
    Baseline & $\tanh$, $4\times64$          & $13.7{\pm}5.4$ & $0.77{\pm}0.23$ & $8.41{\pm}5.41$ & $0.73{\pm}0.20$\\
    Sine     & $\sin$, $4\times64$           & $4.46{\pm}1.88$ & $0.74{\pm}0.19$ & $8.85{\pm}4.27$ & $0.73{\pm}0.18$\\
    \midrule
    \multicolumn{6}{l}{\emph{Exact IC (5 seeds, 3000 epochs)}}\\
    Tanh     & $2\times32$                   & $20.6{\pm}6.1$ & $0.34{\pm}0.11$ & $20.4{\pm}7.7$ & $0.30{\pm}0.11$\\
    Tanh     & $4\times64$                   & $13.1{\pm}4.4$ & $0.35{\pm}0.20$ & $6.40{\pm}2.9$ & $0.19{\pm}0.11$\\
    Tanh     & $6\times64$                   & $11.4{\pm}5.9$ & $0.43{\pm}0.20$ & $6.98{\pm}1.7$ & $0.18{\pm}0.08$\\
    Tanh     & $4\times128$                  & $8.75{\pm}4.1$ & $0.33{\pm}0.15$ & $8.62{\pm}5.1$ & $0.28{\pm}0.16$\\
    Sine     & $2\times32$                   & $5.96{\pm}1.7$ & $0.27{\pm}0.10$ & $10.4{\pm}4.1$ & $0.23{\pm}0.10$\\
    Sine     & $4\times64$                   & $4.85{\pm}2.1$ & $0.27{\pm}0.15$ & $5.80{\pm}1.6$ & $0.21{\pm}0.10$\\
    Sine     & $6\times64$                   & $12.1{\pm}7.8$ & $0.27{\pm}0.11$ & $54{\pm}77$ & $0.35{\pm}0.23$\\
    Sine     & $4\times128$                  & $9.62{\pm}5.1$ & $0.25{\pm}0.11$ & $18.0{\pm}6.6$ & $0.28{\pm}0.09$\\
    \midrule
    \multicolumn{6}{l}{\emph{Finite volume}}\\
    \multicolumn{2}{l}{Lax--Wendroff (limited), $\pi/200$} & $1.92$ & $0.09$ & $2.01$ & $0.07$\\
    \bottomrule
  \end{tabular}
  \caption{PINN and finite-volume comparison at $\nu=0.012$. PINN entries are
    mean$\pm$s.d. over five seeds; predictions are not ensembled. All errors
    are evaluated on the common $201\times101$ space-time grid. The soft-IC
    rows use the two $4\times64$ architectures without the exact-IC branch;
    the exact-IC rows use Eq.~\eqref{eq:hardic}. The Lax--Wendroff row is a
  single deterministic solve.}
  \label{tab:arch}
\end{table}

\section{Discussion and conclusion}
\label{sec:discussion}
For this problem the limited finite-volume scheme remains the more accurate
method: about $2\times10^{-3}$ $L^1$ at $\Delta x=\pi/200$, against $10^{-2}$
for the baseline PINN. The PINN's $L^\infty$ error, however, is dominated by
the discontinuous initial datum rather than by the interior solution, and
imposing the datum exactly removes that contribution at no parameter cost. In
the architecture sweep, exact IC reduces the mean $L^\infty$ to
$0.18$--$0.43$, but the finite-volume scheme remains about $2.5$--$3$ times
more accurate in both norms. It does not close either gap, and at
$\nu=10^{-4}$ the moving front and the outflow dominate again; the reference
$L^\infty$ uncertainty at that value ($4\times10^{-2}$--$10^{-1}$) is also
larger than at $\nu=0.012$.

The lowest mean architecture depends on the target: the $4\times64$ sine
network for linear $L^1$, $4\times128$ sine for linear $L^\infty$, and the
$4\times64$ or $6\times64$ $\tanh$ network for nonlinear $L^\infty$; these
differences are within seed variability. None of these choices changes the
overall conclusion. For this problem the limited scheme remains the better
choice. For a PINN, the exact initial condition is the most effective change
for $L^\infty$ among those tested, whereas the sine activation gives the
largest linear $L^1$ improvement; further gains require resolving the moving
front rather than the initial datum.

Figure~\ref{fig:discussion_profiles} shows final-time profiles for both
problems. The exact-IC PINN tracks the reference and the limited scheme over
most of the domain; the visible differences are at the moving front and near
the outflow, where the residual $L^\infty$ error remains.

\begin{figure}
  \centering
  \includegraphics[width=0.90\textwidth]{\resultsdir/discussion_profiles.png}
  \caption{Final-time profiles at $\nu=0.012$. Linear (left): exact-IC
    $4\times64$ sine network (seed 2). Nonlinear (right): exact-IC
    $4\times64$ $\tanh$ network (seed 0). Solid black: reference; red dashed:
    PINN; dotted blue: limited Lax--Wendroff at $\Delta x=\pi/200$. The PINN
  overlaps the reference except near the moving front and the outflow.}
  \label{fig:discussion_profiles}
\end{figure}

\section*{Acknowledgment}
The code and text were developed with assistance from the DeepSeek V4.1
Flash language model.

\bibliographystyle{plainnat}
\setlength{\bibsep}{0pt}
\bibliography{references}

\end{document}
